\documentclass[10pt,a4paper]{amsart}
\usepackage[margin=19mm]{geometry}
\usepackage{amsmath,amssymb,mathtools}
\usepackage[scr=rsfs,cal=euler]{mathalfa}
\usepackage{xfrac}
\usepackage[unicode,hidelinks,bookmarks=false]{hyperref}
\newcommand{\ie}{i.e.\ }
\newcommand{\cf}{cf.\ }
\newcommand{\resp}{respectively\ }
\newcommand{\Subsection}{\subsection}
\providecommand{\nobreakdash}{\nobreak\hbox{-}}
\setlength{\emergencystretch}{2em}
\multlinegap0pt
\usepackage{bm}
\usepackage{bbm}
\usepackage{dsfont}
\usepackage{xspace}
\usepackage{cancel}
\usepackage{mathtools}
\usepackage{amsthm}
\makeatletter
\@ifundefined{theorem}{\newtheorem{theorem}{Theorem}}{}
\@ifundefined{claim}{\newtheorem{claim}{Claim}}{}
\@ifundefined{lemma}{\newtheorem{lemma}[theorem]{Lemma}}{}
\makeatother
\title[Every 2-Subdivision of a Cubic Graph Is Antimagic]{Every 2-Subdivision of a Cubic Graph Is Antimagic}
\author{Fei-Huang Chang, Teng-Da Chang, Zhishi Pan}
\date{Source: arXiv:2608.11723v1 (CC BY 4.0)}
\begin{document}
\maketitle

\noindent\emph{Source and licence.} Every 2-Subdivision of a Cubic Graph Is Antimagic. Version arXiv:2608.11723v1 (CC BY 4.0), distributed under the Creative Commons Attribution 4.0 International licence, as recorded in the arXiv licence metadata for this exact version and shown on the abstract page https://arxiv.org/abs/2608.11723v1. Full rights notice: licenses/arxiv-2608.11723-CC-BY-4.0.txt.\par\smallskip

\noindent\textit{Editorial scope.} Counted unit: the Lemma whose source label is \texttt{lem:orientation-transfer} in the article (source0.tex lines 569--579, source label \texttt{lem:orientation-transfer}), delivered together with the article's own complete original proof of it (source0.tex lines 581--648, one \texttt{proof} environment of 68 lines). No mathematics is written, repaired or re-derived: statement and proof are the article's own text line for line. Required context retained uncounted: none. Every definition, display and label of those lines is retained unchanged. Omitted from this package and not redistributed: 1--568, 580--580, 649--1097. Nothing inside a retained excerpt is omitted or altered: every retained body is delivered exactly as the article's lines are written, so no omission receipt of any kind accompanies this package. Citations: the retained excerpts carry no citation command at all, so no bibliography entry of the article is transcribed and no reference is redistributed. Reference adaptation: none was needed and none was made; every reference inside the delivered unit resolves inside it, so no unresolved reference or citation is produced. Printed slips kept literal: no printed word, symbol, hypothesis or number of the counted statement or of its proof was altered. Layout and class: the article's own class is \texttt{elsarticle} with packages including amsmath, amssymb, amsthm, mathtools, hyperref; the wrapper is the lane's pinned \texttt{amsart} editorial wrapper at 10pt on A4, which restates the theorem-like environments the retained text needs and adds the article's own macro definitions that the retained text needs (none), with extra wrapper packages (bm, bbm, dsfont, xspace, cancel, mathtools). The delivered numbering is the unit's own. Assets: no retained span contains a figure environment, a graphics-inclusion command, a PDF-page inclusion, a table or a TikZ picture; no image file is copied into this package, the package declares no asset dependency and no image file is redistributed with it. Licence: version arXiv:2608.11723v1 of this article is distributed under the Creative Commons Attribution 4.0 International licence, as recorded in the arXiv licence metadata for this exact version and shown on the abstract page https://arxiv.org/abs/2608.11723v1. A full rights notice is delivered at licenses/arxiv-2608.11723-CC-BY-4.0.txt. Characters: every retained excerpt is pure ASCII, so the delivered engine, XeLaTeX with its default Latin Modern text font, reports no missing character.
\begin{lemma}
\label{lem:orientation-transfer}
Let \(g:E(G)\to[m]\) be a bijection such that every vertex sum induced by
\(g\) occurs at most twice. Then there exist an orientation \(D\) of \(G\)
and a bijection
\[
f:E(S_2(G))\longrightarrow[3m]
\]
such that the sums induced by \(f\) at the original vertices of \(G\) are
pairwise distinct.
\end{lemma}

\begin{proof}
For every vertex sum that occurs twice, pair the two vertices receiving
that sum. Pair the remaining vertices arbitrarily; this is possible
because a cubic graph has even order. Let \(M\) be the resulting
auxiliary perfect matching on \(V(G)\).

\begin{claim}
\label{clm:orientation}
The graph \(G\) has an orientation \(D\) such that
\[
d_D^+(v)\in\{1,2\}
\qquad\text{for every }v\in V(G),
\]
and any two distinct vertices with the same \(g\)-sum have different
outdegrees in \(D\).
\end{claim}

\begin{proof}
The multigraph \(G+M\) is \(4\)-regular. Orient an Euler circuit in each
component and then delete the edges of \(M\). Every vertex has outdegree
\(1\) or \(2\). Moreover, if an edge \(uv\in M\) is oriented from \(u\)
to \(v\), then deleting it leaves
\[
d_D^+(u)=1
\qquad\text{and}\qquad
d_D^+(v)=2.
\]
Hence the two vertices in every prescribed equal-sum pair have different
outdegrees.
\end{proof}

Write
\[
E(G)=\{e_1,e_2,\ldots,e_m\}
\]
so that \(g(e_i)=i\), and orient \(e_i\) from \(u_i\) to \(v_i\)
according to \(D\). In \(S_2(G)\), replace \(e_i\) by the path
\[
u_i x_i y_i v_i
\]
and define
\[
f(x_i y_i)=i,\qquad
f(u_i x_i)=m+2i-1,\qquad
f(y_i v_i)=m+2i.
\]
As \(i\) ranges over \([m]\), the labels \(i\), \(m+2i-1\), and \(m+2i\)
form exactly the set \([3m]\), so \(f\) is a bijection.

For an original vertex \(v\), an edge \(e_i\) directed out of \(v\)
contributes \(m+2i-1\), whereas an edge directed into \(v\) contributes
\(m+2i\). Since \(G\) is cubic,
\[
\sigma_f(v)=3m+2\sigma_g(v)-d_D^+(v).
\]
If two original vertices \(u\) and \(v\) satisfy
\(\sigma_f(u)=\sigma_f(v)\), then
\[
2\bigl(\sigma_g(u)-\sigma_g(v)\bigr)
=
d_D^+(u)-d_D^+(v).
\]
The left-hand side is even, while the right-hand side belongs to
\(\{-1,0,1\}\). Thus both sides are zero. If \(u\ne v\), then \(u\) and
\(v\) are the prescribed pair corresponding to their common \(g\)-sum,
contradicting Claim~\ref{clm:orientation}. Hence \(u=v\), and the
original-vertex sums are pairwise distinct.
\end{proof}

\end{document}
